% !TEX root = ../../main.tex
\subsection{测试环境}

FoodShield 系统的测试环境配置如表~\ref{tab:test_env}所示。所有测试均基于 Flask 内置测试客户端完成，使用独立的临时数据库，不修改演示数据库，确保测试结果的可复现性和数据隔离性。

\begin{table}[H]
\centering
\caption{测试环境配置}
\label{tab:test_env}
\begin{tabular}{ll}
\hline
\textbf{项目} & \textbf{配置} \\
\hline
操作系统 & Arch Linux（WSL2，Kernel 6.18.33.1） \\
CPU & Intel i5-13420H（WSL2 分配 2 核 4 线程，2.611\,GHz） \\
内存 & 7.8\,GiB \\
Python 版本 & 3.13 \\
国密算法库 & gmssl \\
数据库 & SQLite 3 \\
测试框架 & Flask test client（隔离数据库） \\
\hline
\end{tabular}
\end{table}

测试方法分三类：功能测试（正向调用各业务接口验证流程正确性）、安全测试（攻击演示脚本 \texttt{attack\_demo.py} 模拟各类威胁场景）、性能测试（基准脚本在不同数据规模下收集关键路径耗时）。
